\chapter*{Résumé}
\addcontentsline{toc}{chapter}{Résumé}

\vspace{1em}
\noindent\textbf{Résumé :} \\[0.5em]
Ce rapport de stage présente la conception et la réalisation d'un laboratoire virtuel de sécurité réseau basé sur GNS3 et les équipements FortiGate, effectué au sein de l'entreprise Axeli à Rabat, Maroc, sur une durée de cinq semaines (du 2 juillet au 8 août 2026). L'objectif principal était de déployer un environnement d'infrastructure réseau sécurisé, virtualisé et reproductible, capable de simuler les fonctionnalités d'un pare-feu FortiGate. La méthodologie a reposé sur l'utilisation de GNS3, QEMU et des machines virtuelles FortiGate, complétées par des conteneurs Docker pour les services applicatifs. Les travaux ont porté sur la configuration de firewall policies, de NAT (SNAT), du routage inter-LAN via Docker Router, de la high availability (HA) en mode actif-passif, ainsi que des profils de sécurité (antivirus, IPS, Web/DNS filtering, SSL inspection). Le monitoring a été assuré via Grafana, Prometheus et syslog. En raison des contraintes de la licence d'évaluation FortiGate (3 interfaces, 3 politiques, 3 routes, absence de FortiGuard), 7 objectifs sur 16 ont été entièrement atteints, 5 partiellement (profils configurés mais non testables), et 4 n'ont pas été réalisables.

\vspace{1em}
\noindent\textbf{Abstract :} \\[0.5em]
This internship report describes the design and implementation of a virtualized network security lab based on GNS3 and FortiGate appliances, carried out at Axeli, an IT company based in Rabat, Morocco, over a five-week period (July 2 to August 8, 2026). The primary objective was to deploy a secure, virtualized, and reproducible network infrastructure environment capable of simulating FortiGate firewall functionalities. The methodology relied on GNS3, QEMU, and FortiGate virtual machines, supplemented by Docker containers for application services. The work encompassed firewall policy configuration, NAT (SNAT), inter-LAN routing via Docker Router, high availability (HA) in active-passive mode, and security profiles (antivirus, IPS, Web/DNS filtering, SSL inspection). Monitoring was implemented through Grafana, Prometheus, and syslog. Due to FortiGate evaluation license constraints (3 interfaces, 3 policies, 3 routes, no FortiGuard services), 7 out of 16 objectives were fully achieved, 5 partially (profiles configured but untestable), and 4 were not feasible.
